\section{Errata and Critical Corrections}
\label{sec:errata-corrections}

This appendix documents critical corrections made to establish mathematical rigor. These corrections address fundamental issues that, if left uncorrected, would invalidate the proof.

\subsection{Correction 1: Scale Setting}

\begin{correction}[Circular Scale Definition]
\textbf{Original (incorrect):} The lattice spacing was defined as $a(\beta) = 1/\Delta_{\text{lat}}(\beta)$ or $a(\beta) = \xi(\beta)$ where $\xi$ is the correlation length.

\textbf{Problem:} This makes $\Delta_{\text{phys}} = a \cdot \Delta_{\text{lat}} = 1$ by definition, which is circular.

\textbf{Correction:} Define $a(\beta)$ via a \textbf{renormalization condition} independent of spectral quantities:
\begin{itemize}
    \item Gradient flow scale: $a(\beta) = \sqrt{8t_0(\beta)}/c$ where $t_0$ is defined by $\langle t^2 G_{\mu\nu}G^{\mu\nu} \rangle|_{t=t_0} = 0.3$
    \item Or: Schr\"odinger functional coupling at fixed physical size
    \item Or: Small Wilson loop matching
\end{itemize}

See Appendix \ref{sec:rigorous-continuum-scaling} for the complete treatment.
\end{correction}

\subsection{Correction 2: Conditional Tensorization}

\begin{correction}[Missing Interaction Term]
\textbf{Original (incorrect):} Claimed that if $\mu_1$ has LSI($\rho_1$) and all $\mu_x$ have LSI($\rho_2$) uniformly, then $\mu$ has LSI($\min(\rho_1, \rho_2)$).

\textbf{Problem:} This ignores the interaction between $x$ and $y$ variables. The gradient of the conditional expectation includes a covariance term that was set to zero without justification.

\textbf{Correction:} The correct bound is:
\begin{equation}
    \rho \geq \left( \frac{1}{\rho_1} + \frac{1 + C_{\mathrm{int}}}{\rho_2} \right)^{-1}
\end{equation}
where $C_{\mathrm{int}} = \sup_x \sup_{|v|=1} \mathrm{Var}_{\mu_x}(\langle v, \nabla_x H \rangle)$ is the interaction constant.

See Appendix \ref{sec:rigorous-conditional-tensorization} for the complete proof.
\end{correction}

\subsection{Correction 3: Cheeger/Casimir Argument}

\begin{correction}[$\beta$-Independent Gap Claim]
\textbf{Original (incorrect):} Claimed the spectral gap is $\beta$-independent because eigenvalues of the Laplacian on $SU(N)$ are Casimir invariants.

\textbf{Problem:} This applies to the \textbf{unweighted} Laplace-Beltrami operator. The Yang-Mills measure has weight $e^{-\beta S}$, so the relevant operator is a Witten Laplacian whose spectrum depends on $\beta$.

\textbf{Correction:} Use weighted analysis:
\begin{itemize}
    \item Bakry-\'Emery criterion with $\mathrm{Ric}_V = \mathrm{Ric} + \mathrm{Hess}(V)$
    \item Note that $\mathrm{Hess}(S)$ is not positive definite (action is non-convex)
    \item Conclude that alternative methods (Dobrushin, interpolation) are required
\end{itemize}

See Appendix \ref{sec:weighted-spectral} for the complete analysis.
\end{correction}

\subsection{Correction 4: Empirical Inputs}

\begin{correction}[Phenomenological Constants]
\textbf{Original (incorrect):} Used $\sigma_{\text{phys}} = (440 \text{ MeV})^2$ as an input for scale setting and gap bounds.

\textbf{Problem:} This is an empirical value from experiment/lattice simulations, not a mathematical proof.

\textbf{Correction:} 
\begin{enumerate}
    \item Define string tension intrinsically: $\sigma_{\text{phys}} := \lim_{a \to 0} a^{-2} \sigma_{\text{lat}}(\beta(a))$
    \item Prove $\sigma_{\text{lat}}(\beta) > 0$ for all $\beta$ (center symmetry argument)
    \item Prove the scaling relation holds along the continuum trajectory
    \item Do \textbf{not} assume any numerical value
\end{enumerate}
\end{correction}

\subsection{Correction 5: Continuum Limit Scaling}

\begin{correction}[Spectral Permanence]
\textbf{Original (incorrect):} Claimed that if $\Delta_{\text{lat}}(\beta) \geq c > 0$ uniformly in lattice units, the continuum has gap $\geq c$.

\textbf{Problem:} The lattice gap in lattice units diverges as $a \to 0$ if the physical gap is finite. The correct condition involves \textbf{scaled} quantities.

\textbf{Correction:} The correct condition is:
\begin{equation}
    1 - \lambda_1(T_a) \geq m \cdot a
\end{equation}
for some $m > 0$ uniform as $a \to 0$. This gives a physical mass gap $m_a \geq m$.

See Appendix \ref{sec:rigorous-continuum-scaling} for the proper Mosco convergence treatment.
\end{correction}

\subsection{Correction 6: Dimensional Transmutation}

\begin{correction}[The Central Analytical Challenge]
\textbf{Original (status):} Identified as ``Open Problem R.90.1'' but not resolved.

\textbf{Problem:} This is the core of the mass gap problem, not a side issue. The naive bound $\Delta_{\text{lat}} \gtrsim 1/\beta$ combined with $a(\beta) \sim e^{-\beta/(2\beta_0 N)}$ gives $\Delta_{\text{phys}} \to \infty$.

\textbf{Correction:} This must be addressed via one of:
\begin{enumerate}
    \item \textbf{Adjoint interpolation:} Anchor at SYM where the gap is known, deform to pure YM
    \item \textbf{Analyticity + strong coupling:} Prove absence of phase transitions, extend strong coupling result
    \item \textbf{Direct weak coupling analysis:} Extremely difficult, requires non-perturbative methods
\end{enumerate}

The proof cannot simply bypass this by redefining $a(\beta)$ in terms of $\Delta$.
\end{correction}

\subsection{Summary of Required Analysis}

\begin{center}
\begin{tabular}{|l|l|l|}
\hline
\textbf{Issue} & \textbf{Status} & \textbf{Resolution} \\
\hline
Scale setting & Fixed & Gradient flow / renorm condition \\
\hline
Conditional tensorization & Fixed & Interaction-aware theorem \\
\hline
Cheeger/$\beta$-independence & Fixed & Weighted Laplacian analysis \\
\hline
Empirical inputs & Fixed & Intrinsic definitions only \\
\hline
Spectral permanence & Fixed & Correct $O(a)$ scaling \\
\hline
Dimensional transmutation & Via interpolation & Adjoint QCD bridge \\
\hline
\end{tabular}
\end{center}

\subsection{The Honest Structure of the Proof}

After these corrections, the proof has the following structure:

\begin{enumerate}
    \item \textbf{Strong coupling ($\beta < \beta_0$):} Mass gap via cluster expansion (Appendix \ref{ch:fix28_mayer_montroll_radius}).
    
    \item \textbf{SYM anchor ($m_f = 0$):} Mass gap from Witten index + supersymmetry (established).
    
    \item \textbf{Interpolation:} Gap preserved along $m_f: 0 \to \infty$ trajectory (requires Vafa-Witten positivity + analyticity + decoupling).
    
    \item \textbf{Continuum limit:} OS reconstruction + Mosco convergence + correct scaling (Appendix \ref{sec:rigorous-continuum-scaling}).
\end{enumerate}

\begin{remark}[Remaining Challenges]
The interpolation step (3) is where the heaviest analytical work remains. This requires:
\begin{itemize}
    \item Uniform bounds on fermionic determinants
    \item Control of the vacuum structure across the interpolation
    \item Exclusion of phase transitions
\end{itemize}
These are addressed in the main text and supporting appendices.
\end{remark}



